\documentclass[aspectratio=169,11pt]{beamer}

\usetheme{Madrid}
\usecolortheme{default}
\usefonttheme{professionalfonts}
\setbeamertemplate{navigation symbols}{}
\setbeamertemplate{footline}[frame number]

\usepackage{amsmath, amssymb, booktabs}

\title{Empirical Design, Experiments, And Frontier Questions}
\subtitle{From Labor-Market Mechanisms To Applied Research Designs}
\author{Special Topic 7: Labor Market Design, Contracting, and Mechanism Design -- Week 6}
\date{}

\begin{document}

\begin{frame}
  \titlepage
\end{frame}

\begin{frame}{Central question}
\begin{itemize}
    \item How can labor-market design interventions be evaluated credibly?
    \item The capstone object is not a generic estimator.
    \item It is mechanism $\rightarrow$ margin $\rightarrow$ counterfactual $\rightarrow$ equilibrium response $\rightarrow$ welfare.
    \item Students leave with a template for theory-backed applied projects on matching, recruiting, contracts, platforms, and public-sector assignment.
\end{itemize}
\end{frame}

\begin{frame}{Mechanism-first empirical design}
\small
\begin{tabular}{p{0.25\linewidth} p{0.29\linewidth} p{0.32\linewidth}}
\toprule
Mechanism & Labor margin & Empirical object \\
\midrule
Deferred acceptance & ranking and stability & assigned rank, blocking risk, match quality \\
Timing rule & offer pressure and congestion & applications, interviews, acceptances \\
Incentive contract & effort, sorting, multitasking & output, entry, retention, task allocation \\
Platform ranking & attention and visibility & views, bids, hires, wages \\
Public priority rule & staffing and legitimacy & hard-to-staff fill, worker welfare \\
\bottomrule
\end{tabular}
\vspace{0.2cm}
\[
  \text{rule} \rightarrow \text{behavioral margin} \rightarrow \text{labor outcome} \rightarrow \text{welfare object}
\]
\end{frame}

\begin{frame}{Counterfactuals and credible evaluation}
\begin{itemize}
    \item A design estimate is always relative to another rule or institution.
    \item Matching: old algorithm, decentralized offers, modified priorities, or capacity changes.
    \item Recruiting: different timing, application limits, signals, or information.
    \item Contracts: different monitoring technology, wage rule, bonus formula, or probation design.
    \item Platforms: ranking, disclosure, recommendation, wage floor, or fee rule.
    \item Public assignment: different balance of worker preferences, priorities, staffing, and transparency.
\end{itemize}
\end{frame}

\begin{frame}{Experiments, audits, and pilots}
\small
\begin{tabular}{p{0.27\linewidth} p{0.32\linewidth} p{0.29\linewidth}}
\toprule
Design & Best use & Main threat \\
\midrule
Field experiment & direct rule or information change & spillovers and anticipation \\
Platform A/B test & ranking, visibility, wage, or message rule & equilibrium reallocation \\
Audit & screening, discrimination, first-stage response & later match quality unobserved \\
Institutional pilot & real assignment or staffing reform & bundled implementation changes \\
Staggered rollout & administrative or legal reform & adoption and trend confounding \\
\bottomrule
\end{tabular}
\vspace{0.2cm}
\begin{itemize}
    \item The right question is what the intervention moves and which agents can respond.
\end{itemize}
\end{frame}

\begin{frame}{Administrative and platform data}
\begin{itemize}
    \item Rich process data reveal margins that final matches hide.
    \item Applications, rankings, interviews, offers, waitlists, bids, messages, exposure, and acceptances show how the rule works.
    \item Linked outcomes add wages, productivity, retention, promotions, training, exit, and public-service performance.
    \item The key advantage is process visibility, not just sample size.
    \item These data help separate mechanism effects from correlations in realized matches.
\end{itemize}
\end{frame}

\begin{frame}{Structural matching and equilibrium evaluation}
\small
\[
y_i(\rho)=d_i(\rho)+b_i(\rho)+e_i(\rho)
\]
\begin{itemize}
    \item $d_i(\rho)$: direct design effect of the rule.
    \item $b_i(\rho)$: behavioral response by workers, firms, programs, or managers.
    \item $e_i(\rho)$: equilibrium adjustment through congestion, sorting, wages, or capacity.
\end{itemize}
\vspace{0.15cm}
\begin{itemize}
    \item Structural evaluation is useful when welfare depends on unchosen matches or unimplemented mechanisms.
    \item Reduced-form and experimental evidence are strongest for local mechanisms.
    \item Best projects make the two speak to each other.
\end{itemize}
\end{frame}

\begin{frame}{Welfare and portability}
\begin{itemize}
    \item A better metric is not automatically a welfare gain.
    \item A stable match may lower wages; faster recruiting may worsen fit; public staffing gains may reduce worker welfare.
    \item Welfare claims must name whose welfare counts and whether surplus rises or rents shift.
    \item Portability should be argued through the mechanism: congestion, information, incentives, assignment, priority design, or bargaining power.
    \item External validity fails when outside options, legal constraints, wage setting, or equilibrium response differ in central ways.
\end{itemize}
\end{frame}

\begin{frame}{Frontier project map}
\small
\begin{tabular}{p{0.27\linewidth} p{0.35\linewidth} p{0.26\linewidth}}
\toprule
Frontier area & Research question & Data/design \\
\midrule
Within-vacancy design & how are applicants screened under bottlenecks? & application-level admin data \\
Platform wage rules & who bears pricing and transparency changes? & A/B tests, rollouts \\
Internal assignment & do rules affect promotions and retention? & personnel panels, pilots \\
Public staffing & can priorities improve service without exit? & assignment and outcome data \\
AI screening & what margin does automation move? & audit plus process data \\
\bottomrule
\end{tabular}
\end{frame}

\begin{frame}{Research Lab design}
\begin{itemize}
    \item Primary anchor: empirical evaluation of the medical match.
    \item Challenge anchor: deferred acceptance in Army officer labor markets.
    \item Reproduce: compare a decentralized early-offer rule with deferred-acceptance-style assignment in synthetic matching data.
    \item Diagnose: separate preference rank, priority alignment, match quality, shortage staffing, spillovers, and welfare.
    \item Transfer: write a frontier design memo for platforms, recruiting, contracts, teacher placement, public staffing, internal assignment, or AI screening.
\end{itemize}
\end{frame}

\begin{frame}{Student project deliverable}
\small
\begin{tabular}{p{0.27\linewidth} p{0.60\linewidth}}
\toprule
Element & What the memo must state \\
\midrule
Institution & the labor market and rule being studied \\
Mechanism & the design object that moves behavior \\
Labor margin & applications, assignment, wages, effort, retention, or staffing \\
Empirical leverage & randomization, pilot, rollout, cutoff, or model restriction \\
Counterfactual & the alternative rule or institution \\
Equilibrium threat & spillovers, sorting, congestion, or strategic response \\
Welfare object & workers, firms, public agencies, clients, or communities \\
\bottomrule
\end{tabular}
\end{frame}

\begin{frame}{Takeaways}
\begin{itemize}
    \item Labor-market design research starts with the mechanism, not the estimator.
    \item Credible evaluation needs a counterfactual and an observable labor object.
    \item Experiments, audits, pilots, process data, and structural models see different pieces of the same design problem.
    \item Welfare requires equilibrium diagnosis and a clear account of who gains or loses.
    \item The course ends with applied project design: theory-backed, institutionally serious, and empirically disciplined.
\end{itemize}
\end{frame}

\end{document}
